\documentclass[10pt,a4paper]{amsart}
\usepackage[margin=19mm]{geometry}
\usepackage{amsmath,amssymb,mathtools}
\usepackage[scr=rsfs,cal=euler]{mathalfa}
\usepackage{xfrac}
\usepackage[unicode,hidelinks,bookmarks=false]{hyperref}
\newcommand{\ie}{i.e.\ }
\newcommand{\cf}{cf.\ }
\newcommand{\resp}{respectively\ }
\newcommand{\Subsection}{\subsection}
\providecommand{\nobreakdash}{\nobreak\hbox{-}}
\setlength{\emergencystretch}{2em}
\multlinegap0pt
\usepackage{bm}
\usepackage{bbm}
\usepackage{dsfont}
\usepackage{xspace}
\usepackage{cancel}
\usepackage{mathtools}
\usepackage{amsthm}
\makeatletter
\@ifundefined{thm}{\newtheorem{thm}{Thm}}{}
\@ifundefined{lem}{\newtheorem{lem}[thm]{Lemma}}{}
\makeatother
\title[Flat subsets of Euclidean buildings]{Flat subsets of Euclidean buildings}
\author{Raphael Appenzeller, Auguste H\'ebert, Alexander Lytchak}
\date{Source: arXiv:2512.11658v2 (CC BY 4.0)}
\begin{document}
\maketitle

\noindent\emph{Source and licence.} Flat subsets of Euclidean buildings. Version arXiv:2512.11658v2 (CC BY 4.0), distributed under the Creative Commons Attribution 4.0 International licence, as recorded in the arXiv licence metadata for this exact version and shown on the abstract page https://arxiv.org/abs/2512.11658v2. Full rights notice: licenses/arxiv-2512.11658-CC-BY-4.0.txt.\par\smallskip

\noindent\textit{Editorial scope.} Counted unit: the Lem whose source label is \texttt{lem: final} in the article (source0.tex lines 373--380, source label \texttt{lem: final}), delivered together with the article's own complete original proof of it (source0.tex lines 383--398, one \texttt{proof} environment of 16 lines). No mathematics is written, repaired or re-derived: statement and proof are the article's own text line for line. Required context retained uncounted: none. Every definition, display and label of those lines is retained unchanged. Omitted from this package and not redistributed: 1--372, 381--382, 399--490. Nothing inside a retained excerpt is omitted or altered: every retained body is delivered exactly as the article's lines are written, so no omission receipt of any kind accompanies this package. Citations: the retained excerpts carry no citation command at all, so no bibliography entry of the article is transcribed and no reference is redistributed. Reference adaptation: none was needed and none was made; every reference inside the delivered unit resolves inside it, so no unresolved reference or citation is produced. Printed slips kept literal: no printed word, symbol, hypothesis or number of the counted statement or of its proof was altered. Layout and class: the article's own class is \texttt{amsart} with packages including graphics, amssymb, hyperref; the wrapper is the lane's pinned \texttt{amsart} editorial wrapper at 10pt on A4, which restates the theorem-like environments the retained text needs and adds the article's own macro definitions that the retained text needs (none), with extra wrapper packages (bm, bbm, dsfont, xspace, cancel, mathtools). The delivered numbering is the unit's own. Assets: no retained span contains a figure environment, a graphics-inclusion command, a PDF-page inclusion, a table or a TikZ picture; no image file is copied into this package, the package declares no asset dependency and no image file is redistributed with it. Licence: version arXiv:2512.11658v2 of this article is distributed under the Creative Commons Attribution 4.0 International licence, as recorded in the arXiv licence metadata for this exact version and shown on the abstract page https://arxiv.org/abs/2512.11658v2. A full rights notice is delivered at licenses/arxiv-2512.11658-CC-BY-4.0.txt. Characters: every retained excerpt is pure ASCII, so the delivered engine, XeLaTeX with its default Latin Modern text font, reports no missing character.
\begin{lem} \label{lem: final}
	Let $X$ be a CAT(0) space.  Let $Y_1,Y_2 $ be closed, convex subsets of $X$ isometric to convex subsets of some Euclidean spaces.  Let $x\in  Y_1\cap Y_2$ be a point and assume that
	\begin{enumerate} 
	\item $T:={\operatorname{Log}}_x (Y_1) \cup {\operatorname{Log}}_x (Y_2) \subset T_xX \,$ is a convex subset of $T_xX$,
	\item ${\operatorname{Log}}_x (Y_1) \cap {\operatorname{Log}}_x (Y_2) ={\operatorname{Log}}_x (Y_1\cap Y_2)\,.$
	\end{enumerate}
	Then $Y:=Y_1\cup Y_2$ is a  convex subset of $X$ isometric to $T$. 
\end{lem}

	\begin{proof}
		By (2), the logarithmic map restricts to a bijective map
	$$I:=\operatorname{Log}_x :Y\to T\,.$$ 	
		
		
Since $Y_1$ and $Y_2$ are flat and convex, $I$ restricts to isometries
 $$I: Y_i \to I(Y_i)\subset T_x (Y_i)\subset T\subset T_xX\,.$$		 
		
	
		
		Choose arbitrary $y_1,y_2\in Y$ and consider the unique geodesic  $\gamma$ between $I(y_1)$ and $I(y_2)$ in $T$, which exists by (1).  Then $\gamma$ is a concatenation of at most one geodesic in $I(Y_1)$ and at most one geodesic in $I(Y_2)$, hence  the curve $I^{-1}(\gamma)$ has the same length as $\gamma$.  Since $I$ is $1$-Lipschitz, 
		 the curve $I^{-1} (\gamma) \subset Y_1\cup Y_2$  has length $d(y_1,y_2)$.
		 
	 
		 Thus $Y_1\cup Y_2$ is convex and $I:Y\to T$ is an isometry.
	\end{proof}

\end{document}
